% Draft. Uses plain `article` so it compiles anywhere; switch to the ACL/ARR template for submission.
% Every number comes from the repo (scripts/risk2x2.py table, python -m vishell.eval, REPORT.md).
% Red [TODO] marks results that do not exist yet -- do not submit with any left.
\documentclass[11pt]{article}
\usepackage[margin=1in]{geometry}
\usepackage[utf8]{inputenc}
\usepackage[T1,T5]{fontenc}
\usepackage{booktabs}
\usepackage{xcolor}
\usepackage{hyperref}
\usepackage{natbib}
\newcommand{\todo}[1]{\textcolor{red}{[TODO: #1]}}

\title{Asking Is Not Stopping: Small Vietnamese NL-to-Bash Models\\
Learn to Clarify Ambiguity but Not to Avoid Danger}
\author{\todo{author, affiliation}}
\date{}

\begin{document}
\maketitle

\begin{abstract}
A command-line assistant should ask before acting on an underspecified request and stop before running a
harmful one. We ask which of these two behaviours a small language model actually learns. We build a
controlled minimal-pair probe for Vietnamese natural-language-to-Bash: 40 groups of four requests that keep
one operation fixed and vary only whether the request is \emph{clear or vague} and whether its target lies
\emph{inside or outside} an undoable workspace. A 1.5B code model fine-tuned with supervised learning and
with reinforcement learning (GRPO, with and without an undo reward) becomes highly sensitive to vagueness
(ambiguity sensitivity $+0.70$, 95\% CI $[+0.55, +0.82]$) yet remains almost blind to danger: it executes
95\% of clear requests that write to system files, other users' homes or remote hosts (risk sensitivity
$+0.05$, $[0.00, +0.12]$). A large model under a looser protocol reaches $+0.90$. Decoupling the decision
from the generator, with a declarative AST analyzer and a request-only classifier, restores coverage of
dangerous requests; the classifier detects ambiguity well (AUROC 0.89--0.96) but is miscalibrated under
label shift, and checking consistency between intended and actual effects does not detect functionally
wrong commands (AUROC $\le 0.63$). We release the probe, a sandbox-verified Vietnamese benchmark of 158
scenarios, and the analyzer. \todo{update with scale and extension-set results}
\end{abstract}

\section{Introduction}
Small language models are attractive for command-line assistants because they run locally, but a shell
command can destroy data that no later message can recover. Two behaviours matter before execution:
\emph{clarifying} underspecified requests and \emph{stopping} before harmful actions. Standard evaluations
report a single accuracy or refusal rate, which cannot tell the two apart: a model that asks on every vague
request also scores well on many dangerous ones, because dangerous requests are often vague.

We separate them with a minimal-pair design in the spirit of contrast sets~\citep{gardner2020contrast} and
behavioural testing~\citep{ribeiro2020checklist}. Each group fixes one operation (``delete a log file'') and
crosses \emph{clarity} with \emph{risk}: \textit{xoá file build.log} (clear, safe), \textit{xoá file
/var/log/syslog} (clear, risky), \textit{xoá log đi cho nhẹ} (vague, safe), \textit{xoá log của hệ thống cho
nhẹ ổ} (vague, risky). Only the first should be executed directly. Comparing execution rates within a group
isolates how much each factor moves the decision.

Our contributions:
\begin{enumerate}
\item A \textbf{clarity$\times$risk minimal-pair probe} (160 Vietnamese requests, 40 groups) with two
within-group contrasts, risk sensitivity and ambiguity sensitivity, and group-level bootstrap intervals.
\item The finding that \textbf{fine-tuning teaches asking but not stopping}: SFT and GRPO raise ambiguity
sensitivity from $+0.07$ to $+0.70$ while risk sensitivity stays at $+0.05$; an undo-based reward does not
change this. \todo{scale: 0.5B/3B/7B; judge: does the model recognise the risk when shown the command?}
\item An analysis of \textbf{decoupling safety from generation}: a declarative AST analyzer covers dangerous
requests; a request-only classifier ranks ambiguity well but is miscalibrated under label shift; and
intent--effect consistency is \textbf{not} a useful detector of functional errors.
\item A \textbf{sandbox-verified Vietnamese NL-to-Bash benchmark} (158 scenarios, 10--11 paraphrases each,
noisy variants without diacritics or with code-switching, undo commands) and all code.
\end{enumerate}

\section{Related Work}
\paragraph{Natural language to Bash.} NL2Bash~\citep{lin2018nl2bash} introduced a corpus of English
descriptions paired with commands; InterCode~\citep{yang2023intercode} added execution feedback; recent work
revisits evaluation with functional equivalence~\citep{westenfelder2025nl2sh}. These focus on correctness in
English, not on when a system should act at all. \todo{check for existing Vietnamese NL-to-code/Bash resources}
\paragraph{Safety of tool-using agents.} ToolEmu~\citep{ruan2024toolemu} emulates tool execution to surface
risky agent actions; R-Judge~\citep{yuan2024rjudge} benchmarks whether models recognise risk in agent
trajectories. Our judge probe asks a similar recognition question for single commands, but our main object
is the \emph{decision} to execute and its dependence on phrasing versus consequence.
\paragraph{Clarification and ambiguity.} Asking clarifying questions~\citep{aliannejadi2019asking} and
answering ambiguous questions~\citep{min2020ambigqa} study ambiguity in information seeking; we study it
jointly with risk in an action setting.
\paragraph{Shortcuts and controlled evaluation.} Models often exploit surface cues~\citep{gururangan2018annotation,
geirhos2020shortcut}; minimal pairs~\citep{gardner2020contrast} expose such behaviour. Our result is a case of
a model keying on the linguistic form of underspecification rather than on the consequence of the action.

\section{Setup}
\paragraph{Models and training.} All systems start from Qwen2.5-Coder-1.5B-Instruct~\citep{hui2024qwen25coder}.
A: two-stage SFT (translated NL2Bash, then scenario data); B: A followed by GRPO~\citep{shao2024deepseekmath} with a sandbox reward; C: as B plus
a reward for leaving the workspace restorable (undo). Each system answers with one JSON line choosing
\texttt{execute}, \texttt{probe} (read-only) or \texttt{ask}.
\paragraph{Benchmark.} 158 scenario templates (76 execute, 42 probe, 20 ambiguous, 20 irreversible), each with
10--11 Vietnamese paraphrases and noisy variants; every template was certified in a network-less Docker
sandbox (158/158 after two rounds). Training and test splits are disjoint at the template level.
\paragraph{Earlier result: RL adds little.} On held-out templates, GRPO improves execution accuracy over SFT by
$+2.5$ points (95\% CI $[+0.2, +5.5]$) and action accuracy by $+1.3$ (CI includes 0); the undo reward has no
measurable effect. This motivated asking \emph{what} the models learn rather than how to train them harder.

\section{The Clarity$\times$Risk Probe}
Each of 40 groups has one operation and four requests. \emph{Risky} means the target lies outside the
snapshotted workspace (system directories, home, parent directory, other users, devices, git remotes, other
hosts). Items receive opaque shuffled identifiers and are presented one per conversation with the same
instructions (greedy decoding). For a model $m$ and cell $c$, $e_m(c)$ is the fraction of groups executed.
\[
\text{risk sensitivity} = e(\text{clear,safe}) - e(\text{clear,risky}),\qquad
\text{ambiguity sensitivity} = e(\text{clear,safe}) - e(\text{vague,safe}).
\]
Both are 1 for an ideal model and 0 for a model that ignores the factor. Intervals are 95\% percentile
bootstrap over groups (2{,}000 resamples).

\begin{table}[t]
\centering\small
\begin{tabular}{lcccccc}
\toprule
& \multicolumn{4}{c}{execution rate} & \multicolumn{2}{c}{sensitivity [95\% CI]}\\
\cmidrule(lr){2-5}\cmidrule(lr){6-7}
model & clear/safe & clear/risky & vague/safe & vague/risky & risk & ambiguity\\
\midrule
Qwen2.5-Coder-1.5B & 0.97 & 0.97 & 0.90 & 0.80 & $+0.00$ [0.00, 0.00] & $+0.07$ [$-0.03$, 0.18]\\
A: SFT & 1.00 & 0.95 & 0.30 & 0.10 & $+0.05$ [0.00, 0.12] & $+0.70$ [0.55, 0.82]\\
B: SFT+GRPO & 1.00 & 0.95 & 0.30 & 0.10 & $+0.05$ [0.00, 0.12] & $+0.70$ [0.55, 0.82]\\
C: SFT+GRPO+undo & 1.00 & 0.95 & 0.28 & 0.05 & $+0.05$ [0.00, 0.12] & $+0.72$ [0.57, 0.85]\\
\todo{Coder-0.5B / 3B / 7B (4-bit)} & & & & & & \\
Large model$^\dagger$ \todo{name, version} & 0.95 & 0.05 & 0.03 & 0.00 & $+0.90$ [0.80, 0.97] & $+0.92$ [0.83, 1.00]\\
\bottomrule
\end{tabular}
\caption{Clarity$\times$risk probe (40 groups). $^\dagger$Answered in four batches of 40 items within a chat
(one variant per group per batch), not one item per conversation; shown as a reference point only.
\todo{rerun the large model item-by-item through its API}}
\label{tab:probe}
\end{table}

\section{Results}
\paragraph{Fine-tuning teaches asking, not stopping.} Table~\ref{tab:probe}: the base model executes nearly
everything. After SFT, vague requests are executed far less (0.90$\to$0.30 when safe, 0.80$\to$0.10 when
risky), but clear requests are executed regardless of their target (0.95 on clear/risky). GRPO and the undo
reward leave both sensitivities unchanged. Risk has an effect only among vague requests (0.30 vs.\ 0.10),
consistent with danger acting, at most, as an extra cue once a request is already underspecified.
\paragraph{Recognition versus action.} \todo{judge results: when shown the command it would run, does each model
say it leaves the workspace? If yes while still executing, the failure is in acting on knowledge, not in
recognising risk.}
\paragraph{Scale.} \todo{risk sensitivity vs.\ size (0.5B, 1.5B, 3B, 7B, large).}
\paragraph{Independent test set.} \todo{extension set written by other people; inter-annotator agreement
(Cohen's $\kappa$) for clarity and risk; same table.}

\section{Decoupling the Safety Decision}
Given the above, we let the model only generate a command and move the decision to separate components.
\paragraph{Components.} (i) An AST analyzer parses the command with \texttt{bashlex}, walks pipelines,
redirections, command and process substitution, \texttt{bash -c}, \texttt{find -exec}, \texttt{xargs} and
privilege wrappers, and maps each part through a declarative table of $\sim$150 tools to effects (read,
write, overwrite, delete, permission change, process control, network, privilege, remote execution), scopes
(single, glob, recursive, outside-workspace, unresolved) and a risk level; unparseable input, unknown tools and
code executed from text fail closed. (ii) A request-only classifier (XLM-R~\citep{conneau2020xlmr}) predicts
ambiguity and expected effects, trained on labels derived automatically from the AST of gold commands.
(iii) A policy asks when the request is ambiguous, regenerates up to $k$ times on effect mismatches, blocks
unexpected effects that the table itself rates dangerous, and otherwise runs or confirms by risk. Tests
enforce that the final risk equals the analyzer's risk and that the classifier can only make a decision
stricter.
\paragraph{Coverage of dangerous requests.} On the probe, decisions based on the analyzer stop 0.91 of risky
requests for the base model's commands versus 0.07 for the model's own choice; on held-out templates the
figures are 1.00 versus 0.90 for the GRPO model. \textbf{Caveat:} the probe defines risk as leaving the
workspace, which is what the analyzer checks, and two rules were revised after inspecting probe errors; the
analyzer rules are frozen (\texttt{rules-v1}) and \todo{re-evaluate on the extension set}.
\paragraph{Ambiguity is ranked well but miscalibrated.} On the probe the classifier's ambiguity decisions
have precision 1.00 but recall 0.23--0.34 at threshold 0.5, while AUROC is 0.89--0.96. Training data contain
$\sim$13\% ambiguous requests, the probe 50\%; scores are correctly ordered but low (median 0.07 vs.\ 0.005).
This is a prior-shift calibration problem~\citep{saerens2002adjusting}; thresholds should be chosen on
validation data, and we report AUROC. \todo{val-chosen threshold results}
\paragraph{Intent--effect consistency does not detect wrong commands.} Using sandbox execution on held-out
templates as ground truth, the mismatch flag has AUROC 0.36--0.43 for bash-only generations and 0.47--0.54
for JSON-format models (continuous score: 0.27--0.40 and 0.56--0.63). Wrong commands usually have the right
\emph{kind} of effect and the wrong \emph{target} (writing \texttt{/etc/hosts.dev} instead of
\texttt{./hosts.dev}). \todo{pilot: target/argument consistency}

\section{Limitations}
All templates and probe items were written by one author; the extension set and a second annotator are
needed to rule out author-specific phrasing. Only the Qwen2.5-Coder family is tested, and the 7B model is
4-bit quantised. The large-model reference uses a different protocol. The analyzer counts option values as
targets, does not track variable assignments, and fails closed on constructs \texttt{bashlex} cannot parse;
its sandbox backend on Linux isolates the network but not the filesystem.

\section*{Ethics Statement}
Commands generated during evaluation run only inside a network-less sandbox; the analyzer is designed to
fail closed. The probe contains requests to modify system files; it contains no personal data.

\bibliographystyle{plainnat}
\bibliography{refs}
\end{document}
